\documentclass[msom,blindrev]{informs3_ExtAbst}

\OneAndAHalfSpacedXI % current default line spacing
%%\OneAndAHalfSpacedXII 
%%\DoubleSpacedXII
%%\DoubleSpacedXI

% If hyperref is used, dvi-to-ps driver of choice must be declared as
%   an additional option to the \documentstyle. For example
%\documentclass[dvips,mnsc]{informs3}      % if dvips is used
%\documentclass[dvipsone,mnsc]{informs3}   % if dvipsone is used, etc.

% Private macros here (check that there is no clash with the style)

% Natbib setup for author-year style
\usepackage[numbers]{natbib}
 % \bibpunct[, ]{(}{)}{,}{a}{}{,}%
 % \def\bibfont{\small}%
 % \def\bibsep{\smallskipamount}%
 % \def\bibhang{24pt}%
 % \def\newblock{\ }%
 % \def\BIBand{and}%
\usepackage{array}
%% Setup of theorem styles. Outcomment only one.
%% Preferred default is the first option.
\TheoremsNumberedThrough     % Preferred (Theorem 1, Lemma 1, Theorem 2)
%\TheoremsNumberedByChapter  % (Theorem 1.1, Lema 1.1, Theorem 1.2)
\ECRepeatTheorems
\usepackage{bbm}
\usepackage[table,xcdraw]{xcolor}
\usepackage{enumitem}
\usepackage{bm} 
\usepackage{subcaption}
\usepackage{caption}
\usepackage{graphicx}
\newcommand{\NS}[1]{\textcolor{red}{#1}}
%% Setup of the equation numbering system. Outcomment only one.
%% Preferred default is the first option.
\EquationsNumberedThrough    % Default: (1), (2), ...
%\EquationsNumberedBySection % (1.1), (1.2), ...

% For new submissions, leave this number blank.
% For revisions, input the manuscript number assigned by the on-line
% system along with a suffix ".Rx" where x is the revision number.
% \MANUSCRIPTNO{MS-0001-1922.65}


% Packages and Commands
\usepackage[xindy, acronyms]{glossaries}   % Glossaries and 
\usepackage{booktabs}

\usepackage{url}
\def\UrlBreaks{\do\/\do-}
\usepackage{breakurl}
\usepackage[breaklinks]{hyperref}
\newcommand{\E}{\mathbb{E}}
\newcommand{\Pb}{\mathbb{P}}
\usepackage{graphicx} % Required for inserting images
\usepackage{hyperref}

\providecommand{\newblock}{}


% \usepackage{geometry}
% \geometry{a4paper, margin=1in}
% \usepackage{hyperref}
% \usepackage{graphicx}
% \usepackage{xcolor}
% \usepackage{amsmath}
% \usepackage{amssymb}
% \usepackage{algorithm}
% \usepackage{algpseudocode}

% --- Clean Todo Box Setup ---
\usepackage[most]{tcolorbox}

\newtcolorbox{todoblock}{
  colback=gray!8!white,       % Light grey background (change to blue!5!white for soft blue)
  colframe=gray!30!white,     % Subtle border that blends nicely
  boxrule=1pt,                % Thin border
  arc=4pt,                    % Softly rounded corners
  width=\linewidth,           % Forces it to be exactly as wide as the text
  boxsep=8pt,                 % Breathing room inside the box
  left=10pt, right=10pt, top=8pt, bottom=8pt,
  before skip=12pt,           % Spacing before the box
  after skip=12pt             % Spacing after the box
}
% ----------------------------

\begin{document}



\TITLE{Evaluation and Homogeneity in LLMs}
% \vspace{-3em}
\maketitle

\section{Model}
Our core model consists of $N$ firms, each of which competes on $K$ dimensions.
Each firm has a budget $B_{i}$, where they can spend their effort $\mathbf{x}_{i} \in \mathbb{R}^{K}$ across the $K$ dimensions.
Here, putting in effort $x_{i,k}$ produces an output $\theta_{i,k} = f(x_{i,k}).$
The goal is to maximize performance according to a central evaluator, which ranks models based on ``win-rate''; performance against other models.
Here, battles are against other models on the $K$ different topics, each with probability $w_{k}$, where $\mathbf{w} \in \Delta$ is a vector set by the central evaluator based on submitted prompts. 
That is, firm $i$ selects $\mathbf{x}_{i}$ to maximize the Bradley-Terry win-rate against other models: 
\[
\max_{\mathbf{x}_{i}} \sum_{i' \neq i} \mathbb{E}_{y \sim \mathrm{Cat}(\mathbf{w})}[\sigma(f(x_{i,y}) - f(x_{i',y}))].
\]
Here, $\sigma$ is a Sigmoid function capturing Bradley-Terry preferences between voters, and $y$ is a random category (e.g., ``Math'' prompts). 
This is subject to a budget constraint: 
\[
\sum_{k=1}^{K} x_{i,k} \leq B_{i}.
\]
We denote the social welfare according to a user distribution $\mathbf{v} \in \Delta$; for a user $y \sim \mathrm{Cat}(\mathbf{v})$, they receive utility $\max \theta_{i,y}$, so that the social welfare is 
\[
\sum_{k=1}^{K} v_{k} \max\limits_{i}  \theta_{i,k}
\]

\begin{todoblock}
    Currently, we assume that the objective function is only a function of $\mathbf{w}$; it would be interesting to see what would happen if the objective function is some combination of $\mathbf{w}$ to $\mathbf{v}$, and add an extra strategic element.  
\end{todoblock}

\section{Equilibrium under $f(x)=\log x$}
As intuition for the problem, we sketch the solution/equilibrium for this problem in the setting where $f(x) = \log x$. 
Our goal is to find what the equilibrium is for $\mathbf{x}_{i}$; how much effort will each firm put in? 
We first show that the effort allocation factorizes as $x_{i,k} = \alpha_{i} w_{k}$, where $\alpha_{i}$ is firm-specific. 
\begin{lemma} 
The equilibrium solution is $x_{i,k} = \alpha_{i} w_{k}$, where $\alpha_{i}$ is a firm-specific coefficient that increases with larger budget $B_{i}$. 
\end{lemma}
\begin{proof}
    \textbf{Sketch:} 
    Consider the optimization problem for each firm $i$
    \[
    \min -\sum_{i'} \sum_{k} w_{k} \sigma(f(x_{i,k})-f(x_{i',k})),
    \]
    with the budget constraint $\sum x_{i,k} = B_{i}$. 
    Then KKT conditions with $f(x) = \log x$ show that 
    \[
    w_{k} \sum_{i'} \frac{x_{i',k}}{(x_{i,k} + x_{i',k})^2}  = \lambda_{i}-\mu_{i,k}.
    \]
    If $x_{i,k} = \alpha_{i} w_{k}$ (for any $x_{i,k} > 0$), we have that
    \[
    w_{k} \sum_{i'} \frac{\alpha_{i'} w_{k}}{w_{k}^2 (\alpha_{i}+\alpha_{i'})^2} = \sum_{i'} \frac{\alpha_{i'}}{(\alpha_{i}+\alpha_{i'})^2} = \lambda_{i}. 
    \]
    This is constant and only dependent on $i$.
    Since the objective is concave and the constraint set is convex, the KKT conditions are sufficient for a global optimum; the unique optimum is $x_{i,k} = B_{i} w_{k}$, where $\alpha_i = B_i$ follows from the budget constraint $\sum_k x_{i,k} = B_i$ and $\sum_k w_k = 1$.
\end{proof}
\begin{corollary}
    When $f(x) = \log(x)$, firms are maximally homogeneous; $\cos(\mathbf{x}_{i},\mathbf{x}_{j}) = 1$. 
\end{corollary}

Using this structure, we find that budget constrained firms still compete on all dimensions, but are simply Pareto-worse versions of the best model. 
That is, because $\mathbf{w}$ controls the direction of investment, $\alpha_{i}$ only controls the scale of investment. 
Because $\sum x_{i,k} = B_{i}$, we have that $\alpha_{i} = B_{i}$, and so budget-constrained firms simply invest less in each dimension. 

We next show that when $f(x) = \log x$, the welfare loss is proportional to $\mathrm{KL}(\mathbf{v}\|\mathbf{w})$:
\begin{lemma}
    When $f(x) = \log x$, if $\mathbf{x}_{i}$ is the equilibrium with evaluator weights $\mathbf{w}$, and $\mathbf{x}^{*}_{i}$ is the socially optimal allocation (i.e., the equilibrium under $\mathbf{w} = \mathbf{v}$), then the social welfare loss is
    \[
    \sum_{k} v_{k}  \max\limits_{i} f(x^*_{i,k}) - \sum_{k} v_{k} \max\limits_{i} f(x_{i,k}) = \mathrm{KL}(\mathbf{v} \| \mathbf{w}).
    \]
\end{lemma}
\begin{proof}
    \textbf{Sketch:} Because all firms invest in the same directions, the optimal firm in each dimension is firm 1 (assuming $B_{1} \geq B_{2} \cdots B_{N}$), so only firm 1's investment determines welfare.
    The equilibrium gives $x_{1,k} = B_1 w_k$, and the socially optimal allocation is $x^{*}_{1,k} = B_{1} v_{k}$.
    The welfare loss is
    \[
    \sum_{k} v_{k} \bigl[\log(B_1 v_k) - \log(B_1 w_k)\bigr] = \sum_{k} v_{k} \log \frac{v_{k}}{w_{k}} = \mathrm{KL}(\mathbf{v} \| \mathbf{w}).
    \]
\end{proof}

\begin{todoblock}
    It would be interesting to investigate what happens with increasing $B_{i}$ over time with different $f$ and different settings of increasing $B$. Intuitively, increasing $B$ linearly leads to winrates getting near 50\%, as $\theta$ values compress.  
\end{todoblock}
\begin{todoblock}
    It would be interesting to see what incorporating price or brand-related factors would induce as an equilibrium 
\end{todoblock}
\begin{todoblock}
    It would be interesting to see what would happen if $w$ varied over time, or in response to model performance somehow. 
\end{todoblock}

\section{General $f$ Equilibrium}
Computing the explicit equilibrium for general $f$ is messy because $\sigma(f(x_{i,k})-f(x_{j,k}))$ cannot be factored to only depend on $i$ and $j$. 
However, $f$ directly impacts the level of homogeneity through the function $\Omega_{f} = xf'(x)$; when $\Omega_{f}$ is constant, solutions are homogeneous, while when $\Omega_{f}$ varies, solutions are heterogeneous. 
To capture this intuition, we start from the KKT solution, and approximate $c = \frac{\lambda_{i}}{\sum_{i'} \sigma'(f(x_{i,k})-f(x_{i',k}))}$ as a constant: 
\[
f'(x_{i,k}) = \frac{c}{w_{k}}.
\]
Then we can write 
\[
\frac{\mathrm{d}x_{i,k}}{\mathrm{d} \log w_{k}} = \frac{x}{1-\frac{d \log \Omega_{f}}{d \log x}}.
\]
When $\frac{d \log \Omega_{f}}{d \log x} = 0$, $x_{i,k}$ scales with maximum homogeneity (which exactly corresponds to $f(x) = \log x$). 
However, when $\Omega_{f}(x)$ is increasing in $x$, there are higher returns for larger $x$, implying that $x_{i,k}$ is increasingly concentrated at larger $w_{k}$ for higher budgets. 

Formally, we can show comparative statics depending on $\Omega_{f}(x)$: 
\begin{theorem}
    If, for any $a,b$, we have that 
    \[
    |\Omega_{f}(a)-\Omega_{f}(b)| \leq |\Omega_{g}(a)-\Omega_{g}(b)|,
    \]
    then $\mathrm{Var}[\log f'(x_{i,k})-\log f'(x_{j,k})] \leq \mathrm{Var}[\log g'(x_{i,k})-\log g'(x_{j,k})] $.
\end{theorem}
\begin{proof}
    \textbf{Sketch:} 
    By KKT, we can write this as 
    \[
    H_{i,j} = \mathrm{Var}[\log S_{j,k}-\log S_{i,k}] \leq \max_{k,k'} | (\log S_{j,k} - \log S_{i,k}) - (\log S_{j,k'}-\log S_{i,k'})|. 
    \]
    By the Lipschitzness of $\sigma'$, and noting that $\sigma' > S > 0$, we can bound this as 
    \begin{align} 
    \max_{k,k'} | (\log S_{j,k} - \log S_{i,k}) - (\log S_{j,k'}-\log S_{i,k'})| \leq \frac{L}{S}|(f(x_{i,k})-f(x_{j,k})) - (f(x_{i,k'})-f(x_{j,k'}))| \\ 
    \leq \frac{L}{S} \sup\limits_{a,b} | \Omega_{f}(a)-\Omega_{f}(b)| \cdot | \log x_{i,k}-\log x_{i,k'}|
    \end{align}
    The last depends only on $f$ through $|\Omega_{f}(a)-\Omega_{f}(b)|$. 
    Note that the full proof requires showing tightness in this bound; currently we just show that the upper bound is larger, but a full proof needs to show that this implies differences in $H_{i,j}$ itself. 
\end{proof}

We can similarly show that the optimal $w$ for welfare depends on $\Omega$; because the value of $\Omega$ induces under or overinvestment in certain $w_{k}$, inverting $\Omega$ and setting $w$ can undo this change. 
We can show this formally: 
\begin{theorem}
    The optimal $w^{*}_{k} \propto \frac{v_{k}}{\Omega_{f}(x^{*}_{k})/\bar{\Omega}}$, where
    \[
    \bar{\Omega} = \frac{1}{K} \sum_{k=1}^{K} v_{k} \Omega_{f}(x^{*}_{k}).
    \]
\end{theorem}
\begin{proof}
    \textbf{Sketch:} The welfare $W$ is 
    \[
    W = \sum_{k} v_{k} \max_{i} f(x_{i,k}(\mathbf{w})).
    \]
    We can find the optimal $\mathbf{w}$ by Taylor approximating $\frac{\mathrm{d}W}{\mathrm{d}w_k}$ as
    \[
    \frac{\mathrm{d}W}{\mathrm{d}w_k} = v_{k} f'(x^{*}_{k}) \frac{\mathrm{d} x^{*}_{k}}{\mathrm{d} \log w_{k}}.
    \]
    Setting this equal across all $k$ (simplex first-order condition) and simplifying gives
    \[
    v_{k} \Omega_{f}(x^{*}_{k}) = \frac{1}{K} \sum_{k'} v_{k'} \Omega_{f}(x^{*}_{k'}) = \bar{\Omega}.
    \]
    From the first-order approximation $f'(x^*_k) = c/w_k$ we have $w_k \propto 1/f'(x^*_k) = x^*_k/\Omega_f(x^*_k)$.  Substituting $\Omega_f(x^*_k) = \bar{\Omega}/v_k$ (from above) gives
    \[
    w_{k} \propto \frac{x^{*}_{k}}{\Omega_{f}(x^{*}_{k})} = \frac{v_k \, \Omega^{-1}_{f}(\bar{\Omega}/v_k)}{\bar{\Omega}}.
    \]
\end{proof}

\section{Empirical Verification}
We run experiments with the Arena dataset to understand what this model looks like in practice.

\paragraph{Win Rate Decay}
We first investigate win-rate decay and what drives it.
For each landmark frontier model, the ELO advantage above the field median decays approximately linearly over time.
This is consistent with the model under $f(x) = \log x$ and exponentially growing $B_i$: if each firm's budget grows at the same exponential rate, the ELO gap above median falls linearly.
The compression is substantial---win-rate variance among the top-20 Arena models has fallen approximately $6\times$ from 2023 to 2026 (ELO standard deviation roughly halved)---and is symmetric across organizations, consistent with genuine cross-lab convergence in budgets rather than any single lab's strategic behavior.

\paragraph{Production function $f$}
Chinchilla scaling laws~\citep{chinchilla} predict that benchmark performance grows as $\log(\text{compute})$, which implies $f(x) = \log x$ in our model.
This is empirically supported: Arena ELO for the best model grows linearly in time (equivalent to $\log$ in compute if compute scales exponentially), and benchmark scores fit a $\log$ curve in parameter count with $r^2 \approx 0.7$--$0.9$.
Under $f(x) = \log x$, Lemma~1 predicts maximal homogeneity ($\cos(\mathbf{x}_i, \mathbf{x}_j) = 1$).
This is borne out: pairwise Spearman correlation of BBH subtask accuracy vectors across frontier models from different organizations is 0.848 ($p < 0.0001$), and cross-category win-rate rank correlations in Arena battles range from 0.787 to 0.910---consistent with a largely one-dimensional capability surface.

\paragraph{Differences in $w$ and $v$}
Under $f(x) = \log x$, Lemma~2 gives $w^{*} = v$ with welfare loss $\mathrm{KL}(\mathbf{v} \| \mathbf{w})$.
We estimate $\mathbf{w}$ from Arena leaderboard vote counts by category and $\mathbf{v}$ from the distribution of task types in 116,000 organic Arena battles.
The two distributions differ most in math: $\mathbf{w}$ assigns roughly $1.5\times$ more weight to math than $\mathbf{v}$, so the recommendation is to downweight math in the evaluation to reduce the $\mathrm{KL}$ welfare loss.
\NS{NOTE: We need to do a bit more research to understand whether this discrepancy is actually the best estimate of $\mathbf{w}$ and $\mathbf{v}$, or if we can find a better estimate.} 

\section{Related Work}
\paragraph{Monoculture in Markets}
Prior work investigates monoculture in markets, including for matching and hiring~\citep{monoculture_matching,monoculture_social_welfare}. 
Our work similarly studies monoculture in an LLM market, but differs in that we study how \textit{evaluation} impacts these. 

\paragraph{Markets for LLMs}
Work using economics tools to analyze LLMs have looked at the relationship between LLM competition and social welfare~\citep{bayes_risk}, competition in marketplaces~\citep{competition_alignment}, and how model aggregations impact diversity~\citep{aggregation_model_diversity}. 
Our work differs in that we analyze the impact of evaluation as an information signal, with competition rather than aggregation between models. 

\paragraph{Blotto Games}
Blotto games study how different firms should allocate resources under competition, and includes Lottery Blotto, which mimics the same probabilistic win rate~\citep{lottery_blotto}; our key difference is that we study how the evaluation mechanism (the choice of $\mathbf{w}$) shapes equilibrium allocations and social welfare.
Lottery Blotto also implicitly focuses only on the $f(x) = \log x$ case. 

\paragraph{Pricing and LLMs}
Prior work on pricing in language models examines how providers should offer contracts, and shows that contracts should be per-token rather than user-specific~\citep{menu_pricing}. 
Our work differs in our focus on firm allocations rather than user behavior, and on evaluation rather than pricing. 

\bibliographystyle{apalike}
\bibliography{references}

\end{document}

